\providecommand{\vf}[1]{\textcolor{blue}{#1}}
\chapter{Introduction}\label{ch:intro}

A client that sends its prompts to a language model in the cloud exposes their text to the service provider.
To keep that text on the client, split inference partitions the model between the client and the server~\cite{he2019model,mai2024split,nguyen2026noir}.
The client runs the first blocks of the model on its own device and sends the server the hidden states that those blocks produce, one vector per token.
The server runs the remaining blocks on those hidden states.
In some designs the client also runs and fine-tunes the last blocks, so that the server sees neither the prompt nor the generated text~\cite{nguyen2026noir}.
In each of these designs, the server receives hidden states in place of tokens.

However, what split inference protects depends on what the transmitted hidden states reveal.
The server, which is also the adversary, is honest-but-curious: it follows the protocol and still tries to learn what the client wrote.
By observing the hidden states, the server aims to reconstruct the prompt that produced them, an attack known as inversion.
Thus, inversion poses a privacy risk to split inference, and that risk depends on how difficult inversion is and on what the adversary must hold to carry it out.

\paragraph{Prior attacks.} Existing inversion attacks succeed when the adversary has access to the model that produced the hidden states.
With the weights of that model, the adversary optimizes a candidate prompt until its hidden states match the observed ones, or tests vocabulary tokens one position at a time~\cite{qu2025prompt,dong2025depth,thomas2025hidden,nikolaou2026language}.
With query access, the adversary submits its own text, collects the resulting hidden states, and trains an inverter on those pairs~\cite{luo2025prompt,dong2025depth}.

A second line of work inverts sentence embeddings, one pooled vector per text, and again relies on queries to the encoder~\cite{song2020information,li2023sentence,morris2023text}.
Recent attacks on sentence embeddings remove the queries: they align the embeddings of the target encoder with those of an encoder the adversary holds, from thousands of leaked pairs~\cite{huang2024transferable}, from a few~\cite{chen2025algen}, or from none~\cite{jha2025harnessing}.
Those attacks apply to pooled vectors.
To the author's knowledge, no attack of that kind has been evaluated on the per-token hidden states of a language model, the quantity that a split inference client transmits.

\paragraph{Research question.} This dissertation studies an adversary with neither form of access: it holds no weights of the client's blocks and has no means to query them.
The question is how much of a prompt such an adversary reconstructs from the hidden states alone, given a small leak of text.
Because prior attacks rest on access, it is tempting to conclude that a client that keeps its blocks private also keeps its prompts private.
Yet the Platonic Representation Hypothesis proposes that models trained on similar data converge toward similar representations~\cite{huh2024platonic}.
The approach draws on that hypothesis: if it holds for language models, a simple map may relate the hidden states of a target model to those of public models.
An adversary could then build an inverter entirely from public models and connect the target model to it afterwards.
Similar representations do not by themselves expose a prompt: a map can preserve the meaning of a prompt and still lose its exact words.
Prompt reconstruction is therefore a direct test of what survives the map between models.
The dissertation argues that representations shared across models support inversion across models, and that unpaired matching improves reconstruction when a small leak supplies the correspondences.

\paragraph{Approach.} The dissertation develops that attack in two steps.
The first step builds the inverter, a translation hub (Chapter~\ref{ch:hub}).
The hub maps the hidden states of three open-weight language models into one shared space, and each of the three models, kept frozen, reconstructs the prompt from the hub outputs of any other.
The second step connects a target model that took no part in hub training (Chapter~\ref{ch:attack}).
A single affine map sends the hidden states of the target model into the hub, and the frozen decoders of the hub reconstruct its prompts.
Two sources train the map: a small leak, a set of short text spans for which the adversary holds both the text and the hidden states, and a matching term that aligns unpaired hidden states with unpaired hub outputs.
The leak is counted in rows: each row pairs the average hidden state of the target model with the average hub output for one token group of a leaked span.
The matching term requires no text from the client.
A map trained on the leak alone is the supervised fit, and a map trained on the leak and the matching term is the hybrid fit.

\paragraph{Results.} Both steps are evaluated on NQ, the query strings of Natural Questions, and on two datasets of legal text.
The score is word-F1, the overlap between the words of the reconstruction and the words of the prompt, without regard to their order.

The hub supports reconstruction across models: each of its three frozen decoders reconstructs prompts from the hidden states of the other models with high accuracy, on all three datasets.
For the target model, the leak alone carries only a small part of the recovery under supervised training.
The matching term more than doubles that recovery on all three datasets.
Without a leak, the matching term alone reconstructs no better than a control that feeds the decoder the hidden states of another prompt.
Exact match remains the challenging part: the hybrid fit reconstructs few prompts word for word on NQ and nearly none on the legal datasets.

\paragraph{Contributions.}
\begin{itemize}
\item \textbf{An inversion framework centered on a translation hub.} The framework maps the hidden states of several language models into one shared space, the hub.
From the hub, each model, kept frozen, reconstructs the prompts of any other model.
That reconstruction extends to a model outside the hub: a map fitted on a small leak and on unpaired hidden states sends its hidden states into the hub, with no weights of the model and no queries to it.
\item \textbf{A measurement of what an adversary recovers.} This dissertation measures what part of a prompt the framework recovers under different leak scenarios and under structure matching, on short and long prompts.
The proposed work extends that characterization to the privacy methods that defend split inference.
\end{itemize}

\paragraph{Outline.} Chapter~\ref{ch:lit} reviews inversion attacks, cross-model alignment, and defenses, and positions the dissertation among them.
Chapter~\ref{ch:hub} presents the translation hub.
Chapter~\ref{ch:attack} presents the attack on a target model outside the hub.
